\documentclass[11pt]{article}
\usepackage[margin=1in]{geometry}
\usepackage{amsmath,amssymb,mathtools}
\usepackage{fontspec,unicode-math}
\setmainfont{Latin Modern Roman}
\setsansfont{Latin Modern Sans}
\setmathfont{Latin Modern Math}
\usepackage{xurl,hyperref}
\hypersetup{colorlinks=true,urlcolor=blue,linkcolor=blue}
\newcommand{\sref}[2]{\hyperlink{source:#1:#2}{[S#2]}}
\newcommand{\cataloguescope}[1]{\paragraph{Recorded mathematical question.}#1}
\begin{document}
% UnsolvedMath ID 9700018; source catalogue code AMR-096-0018
\setcounter{section}{9700017}
\section{Stretch-length exponent in spatial networks}\label{9700018}
{\small\sffamily UnsolvedMath ID 9700018 · Probability}
\subsection{Definitions and mathematical statement}

\paragraph{Shared notation.}$\mathbb N=\{1,2,\ldots\}$, and $\mathbb Z,\mathbb Q,\mathbb R,\mathbb C$ denote the integers, rationals, reals and complex numbers. Any other conventions are specified locally.


For $n$ distinct cities $z=(z_1,\ldots,z_n)$ in a square of area $n$, a Steiner network $N$ is a connected network of line segments joining them, with additional junctions allowed and crossings usable as junctions. Let $r_N(z_i,z_j)$ be shortest route length and
\[S(N)=\max_{i\ne j}\frac{r_N(z_i,z_j)}{\|z_i-z_j\|_2},\qquad L(N)=\frac{\operatorname{len}(N)}n,\qquad \psi_n(z,s)=\inf_{S(N)\le s}L(N).\]
Network length counts each segment once. For $Z_n$ consisting of $n$ independent uniform points in the area-$n$ square, define the established limits
\[\Psi^{\rm ave}(s)=\lim_{n\to\infty}\mathbb E\psi_n(Z_n,s),\qquad\Psi^{\rm worst}(s)=\lim_{n\to\infty}\sup_z\psi_n(z,s)\quad(s>1).\]
For positive functions $a,b$, $a\asymp b$ means their ratio stays between two fixed positive constants in the indicated limit.
\paragraph{Statement recorded in the supplied source.}
Improve explicit upper and lower bounds for both functions. Is there an exponent $\alpha$ such that $\Psi(s)\asymp(s-1)^{-\alpha}$ as $s\downarrow1$ for each of $\Psi^{\rm ave}$ and $\Psi^{\rm worst}$? Determine whether a common exponent works for both settings. \sref{9700018}{2}
\subsection{Short English statement}
Find sharper optimal network-length bounds under a stretch constraint and identify their divergence exponent as stretch approaches one.
\subsection{Sources}
\begin{description}
\item[\hypertarget{source:9700018:1}{S1}] ULAM.AI and the UnsolvedMath Contributors, \emph{UnsolvedMath: A Curated Collection of Open Mathematics Problems}, record 9700018 (AMR-096-0018). CC BY 4.0. \url{https://huggingface.co/datasets/ulamai/UnsolvedMath}.
\item[\hypertarget{source:9700018:2}{S2}] David Aldous and Tamar Lando, \emph{The stretch-length tradeoff in geometric networks: average case and worst case study}, Mathematical Proceedings of the Cambridge Philosophical Society 159 (2015), 125--145, Section 1, definitions and scaling question, pp. 1--3; Section 2, existence of limits. \url{https://www.stat.berkeley.edu/~aldous/Papers/me142.pdf}.
\end{description}
\label{end:9700018}
\end{document}
